%==============================================================================
%  CAPÍTULO — NORMATIVA EN TRÁMITE
%==============================================================================
\chapter{Normativa en trámite: el derecho concursal que viene}
\label{cap:normativa-tramite}

\fichaCapitulo{Anticipar el cambio normativo.}{La consulta pública como mecanismo de
producción normativa de la Superintendencia, estado del proceso al cierre de esta edición
y criterios para anticipar su impacto en la práctica.}

%==============================================================================
\bloqueTeorico
%==============================================================================

\subsection{Por qué un capítulo sobre normas que aún no existen}

Un tratado que sólo describa el derecho vigente envejece desde el día de su publicación.
En materia concursal ese envejecimiento es especialmente rápido: la \superir{} produce
normativa de forma continua, y una parte relevante de la práctica cotidiana ---formularios,
plazos operativos, deberes de información--- depende de instrumentos que se modifican
varias veces al año.

Existe, además, una razón profesional. La Superintendencia somete sus proyectos
normativos a consulta pública antes de dictarlos. Ese procedimiento cumple dos funciones
para el abogado concursalista: permite anticipar el cambio y, sobre todo, permite
participar en él. Un estudio que tramita concursos con regularidad conoce las
disfunciones operativas mejor que el regulador, y la consulta pública es el cauce
institucional para plantearlas.

\begin{foco}[La consulta pública como herramienta profesional]
Observar un proyecto de norma cuesta un escrito. Litigar durante años contra una norma
mal diseñada cuesta mucho más. El seguimiento de la sección de normativa en trámite
debería formar parte de la rutina del estudio, con la misma regularidad con que se revisa
el \boletin{}.
\end{foco}

\subsection{Naturaleza del trámite de consulta}

La consulta pública no es un procedimiento contradictorio ni vinculante: la
Superintendencia no está obligada a acoger las observaciones, ni a fundamentar su
rechazo en un acto formal. Su valor es distinto: incorpora información que el regulador
no posee y reduce el riesgo de normas que resulten inaplicables en la práctica.

%==============================================================================
\bloqueNormativo
%==============================================================================

\subsection{Estado del proceso al cierre de esta edición}

El siguiente cuadro recoge los proyectos sometidos a consulta pública, con indicación de
su estado \parencite{superir-normativa-tramite}. La columna «estado» distingue entre las
normas cuyo texto definitivo ya fue dictado y las que permanecen en tramitación o fueron
abandonadas.

\begin{table}[htbp]
\centering
\small
\caption{Normativa sometida a consulta pública}
\label{tab:normativa-tramite}
\begin{tabularx}{\textwidth}{@{}p{0.11\textwidth} p{0.15\textwidth} X p{0.11\textwidth}@{}}
\toprule
\textbf{Fecha} & \textbf{Tipo} & \textbf{Materia} & \textbf{Estado} \\
\midrule
03.07.2026 & Instructivo & Presentación de excusas ante la Superintendencia por sujetos
fiscalizados; deroga normas anteriores & Vigente \\
\addlinespace
16.06.2026 & Acto administrativo & Califica como grave la infracción que indica & En
trámite \\
\addlinespace
07.11.2025 & Resolución & Complementa el instructivo sobre prelación de créditos y
repartos de fondos en quiebras & En trámite \\
\addlinespace
07.11.2025 & Resolución & Aprueba instructivo sobre síndicos que cesan anticipadamente
& Dictada \\
\addlinespace
06.10.2025 & NCG & Exigencias del procedimiento de renegociación, acreditación de
información y celebración de audiencias & Dictada como NCG 28 \\
\addlinespace
19.08.2025 & Manual & Recomendaciones y pautas de buenas prácticas para entes
fiscalizados & En trámite \\
\addlinespace
27.11.2024 & Instructivo & Fianzas, preferencias y contragarantías en liquidación con
sociedades anónimas de garantía recíproca & Dictado \\
\addlinespace
20.11.2024 & Resolución & Complementa instructivo sobre notas de débito en prelación y
pago de créditos & Dictada \\
\addlinespace
15.10.2024 & Instructivo & Finiquitos laborales electrónicos y firma electrónica
& Dictado como Instructivo 6 \\
\addlinespace
10.07.2024 & Instructivo & Apertura de cuentas corrientes bancarias solicitadas por
liquidadores & Dictado \\
\addlinespace
10.07.2024 & Instructivo & Incautación, conservación y entrega de bienes & Dictado como
Instructivo 5 \\
\addlinespace
02.07.2024 & Res. exenta & Forma y contenidos de las cuentas provisorias de liquidadores
& Dictada como NCG 27 \\
\addlinespace
02.07.2024 & Resolución & Actuaciones de liquidadores ante el Registro Nacional de
Deudores de Pensiones de Alimentos & Dictada \\
\addlinespace
02.07.2024 & Instructivo & Instrucciones a veedores en reorganización judicial,
simplificada y extrajudicial & Dictado como Instructivo 3 \\
\addlinespace
24.06.2024 & Res. exenta & Instrucciones a síndicos sobre sobreseimiento definitivo y
temporal & Dictada como Instructivo 2 \\
\addlinespace
24.06.2024 & Instructivo & Hechos relevantes que deben informar los sujetos fiscalizados
& Dictado como Instructivo 4 \\
\bottomrule
\end{tabularx}
\end{table}

\verificar{Actualizar esta tabla inmediatamente antes de la impresión: es el contenido de
la obra con mayor velocidad de obsolescencia. Confirmar además la correspondencia entre
cada proyecto y la norma finalmente dictada.}

\subsection{Lectura del cuadro: qué revela sobre la dirección del regulador}

Tres tendencias se desprenden de la secuencia de proyectos:

\begin{enumerate}
  \item \textbf{Digitalización de las actuaciones.} Finiquitos electrónicos, juntas
  telemáticas, plataformas de enajenación y notificación electrónica apuntan a un
  procedimiento íntegramente desmaterializado.

  \item \textbf{Intensificación del control sobre los auxiliares.} Hechos relevantes,
  indicadores de desempeño, cuentas provisorias y régimen de excusas configuran una
  fiscalización más densa y basada en información periódica antes que en revisión
  \emph{ex post}.

  \item \textbf{Articulación con otros registros públicos.} Las instrucciones sobre el
  Registro Nacional de Deudores de Pensiones de Alimentos muestran la creciente
  interconexión del procedimiento concursal con registros ajenos al sistema.
\end{enumerate}

%==============================================================================
\bloquePractico
%==============================================================================

\subsection{Rutina de seguimiento normativo}

\begin{checklist}[Control mensual sugerido]
  \item Revisar la sección de normativa en trámite del sitio de la \superir{}.
  \item Identificar los proyectos que inciden en los procedimientos que el estudio tramita.
  \item Anotar el plazo para presentar observaciones.
  \item Contrastar el proyecto con las disfunciones operativas registradas en las causas
  propias.
  \item Presentar observaciones cuando el proyecto afecte plazos, formularios o deberes
  de información.
  \item Verificar la publicación de la norma definitiva y su fecha de entrada en vigencia.
  \item Actualizar los modelos y listas de verificación del estudio.
\end{checklist}

\subsection{Cómo redactar observaciones útiles}

\begin{practica}[Qué persuade al regulador]
Las observaciones que prosperan no son las que invocan principios generales, sino las que
describen un problema operativo concreto, cuantifican su efecto y proponen una redacción
alternativa. «La exigencia del literal c) obliga a obtener un certificado que el Registro
Civil emite en quince días hábiles, plazo superior al que la propia norma concede para
presentar» es un argumento que se atiende. «La norma vulnera el debido proceso» no lo es.
\end{practica}

\subsection{Gestión del cambio normativo en el estudio}

\begin{checklist}[Cuando se dicta una norma nueva]
  \item Determinar la fecha de entrada en vigencia y el régimen transitorio.
  \item Identificar las causas en tramitación que resultan afectadas.
  \item Actualizar los modelos de escritos y las listas de verificación.
  \item Comunicar el cambio a los clientes cuyas causas se vean alteradas en plazos o
  exigencias.
  \item Registrar la fecha de verificación en el repositorio de normativa del estudio.
\end{checklist}
